\begin{notechapter}
\chapter{Output Distributions and their Losses}\label{ch:output-distributions}
\chaptersubtitle{Choose a distribution, read off the loss}

\begin{chapterintro}
\Cref{ch:cross-entropy} showed that maximum likelihood minimizes a cross-entropy.
Turned around, this gives a recipe for designing losses: decide which family of
distributions the target should follow given the input, let the network output
the parameters of that distribution, and use the negative log-likelihood as the
loss. The squared error, the absolute error, the binary and categorical
cross-entropies and the Poisson loss are all instances. You will learn
\begin{itemize}[nosep]
  \item the recipe, and the role of the \emph{output unit} that maps unconstrained
        network outputs to valid parameters;
  \item the losses of the Gaussian (with fixed or learned variance), Laplace,
        Bernoulli, categorical and Poisson models, with their gradients;
  \item why all of these gradients have the form ``prediction minus target'': the
        exponential family and canonical links;
  \item mixture density networks for multimodal targets, and a word on structured
        outputs.
\end{itemize}
\prerequisites{\cref{ch:cross-entropy}; basic probability distributions.}
\end{chapterintro}

\section{The recipe}\label{output-distributions:sec:recipe}

A network computes from the input $x$ a vector of real numbers $z = g_\theta(x)$,
unconstrained in sign and size. A probabilistic model of $Y$ given $X = x$ needs
parameters with constraints: a probability lies in $(0,1)$, a variance is
positive, the probabilities of $K$ classes sum to $1$. The \emph{output unit} is a
fixed map $\omega = u(z)$ from the network's outputs to the parameters $\omega$ of the
chosen distribution $p(y\mid\omega)$. The loss of one example is the negative
log-likelihood
\[
  \ell(z, y) = -\log p\bigl(y\mid u(z)\bigr),
\]
and training minimizes its average over the data, which by
\cref{cross-entropy:thm:mle} is maximum likelihood. Three choices are involved: the
family, the output unit, and the network. \Cref{output-distributions:tab:summary}
at the end of the chapter collects the usual combinations.

\section{Real-valued targets}\label{output-distributions:sec:real}

\paragraph{Gaussian, fixed variance.} With $p(y\mid\mu) = \mathcal N(y;\mu,\sigma^2)$, $\sigma$ fixed, and
the identity output unit $\mu = z$,
\[
  \ell(z, y) = \frac{(y - z)^2}{2\sigma^2} + \frac12\log(2\pi\sigma^2),\qquad
  \frac{\partial\ell}{\partial z} = \frac{z - y}{\sigma^2}.
\]
Up to an additive constant and the factor $1/\sigma^2$, this is half the squared error, and
its average is half the mean squared error: least squares is maximum likelihood
under Gaussian noise of constant variance. The value of $\sigma$ does not change the
minimizer in $\mu$, and once $\mu$ is fitted the likelihood is maximized by
$\hat\sigma^2 = \frac1n\sum_i(y_i - \hat\mu(x_i))^2$, the mean squared residual
(\cref{output-distributions:exr:1}).

\paragraph{Gaussian, learned variance.} When the noise level depends on $x$
(\emph{heteroscedasticity}), the network outputs two numbers, $z = (z_1, z_2)$, with
$\mu = z_1$ and $\sigma^2 = e^{z_2}$; the output unit $\exp$ guarantees a positive variance,
and $z_2 = \log\sigma^2$ is unconstrained. The loss and its gradients are
\[
  \ell = \frac12z_2 + \frac{(y - z_1)^2}{2e^{z_2}} + \frac12\log2\pi,\qquad
  \frac{\partial\ell}{\partial z_1} = \frac{z_1 - y}{\sigma^2},\qquad
  \frac{\partial\ell}{\partial z_2} = \frac12\Bigl(1 - \frac{(y - \mu)^2}{\sigma^2}\Bigr).
\]
The second gradient vanishes when the predicted variance equals the squared
residual: the network learns to predict its own error. Points with a large
predicted variance contribute less to the gradient of the mean, which makes the
fit robust to regions of high noise.

\begin{pitfall}[The variance can collapse]
On a single example the loss is unbounded below: if $\mu = y$ exactly, letting
$\sigma^2\to0$ sends $\ell\to-\infty$. A flexible network can fit some training points
exactly and shrink their variance without limit. Remedies are a lower bound on
$\sigma^2$, weight decay, fitting the mean first, or a prior on the variance. With a
single variance shared by all points the problem disappears as long as the
residuals are not all zero.
\end{pitfall}

\paragraph{Laplace.} With $p(y\mid\mu) = \frac1{2b}e^{-\abs{y - \mu}/b}$, $b$ fixed, and $\mu = z$,
\[
  \ell(z, y) = \frac{\abs{y - z}}b + \log(2b),\qquad
  \frac{\partial\ell}{\partial z} = \frac{\sign(z - y)}b\quad(z\neq y).
\]
Maximum likelihood is least absolute deviations. The best constant prediction
is a median of the targets (\cref{output-distributions:exr:2},
\cref{bayes-decision-theory:sec:other}), which is insensitive to outliers: the gradient
of each point has the same size however far away it is, whereas under the
Gaussian it grows with the residual. The Laplace density has heavier tails than
the Gaussian, so a large residual is less surprising under it and costs less.

\section{Discrete targets}\label{output-distributions:sec:discrete}

\paragraph{Bernoulli.} For $y\in\{0, 1\}$, $p(y\mid\pi) = \pi^y(1 - \pi)^{1-y}$ with
$\pi = \sigma(z) = 1/(1 + e^{-z})$, the sigmoid. The loss is the binary cross-entropy
\[
  \ell(z, y) = -y\log\sigma(z) - (1 - y)\log\bigl(1 - \sigma(z)\bigr) = \log\bigl(1 + e^{z}\bigr) - yz,\qquad
  \frac{\partial\ell}{\partial z} = \sigma(z) - y .
\]

\paragraph{Categorical.} For $y\in\{1,\dots,K\}$, $p(y\mid\pi) = \pi_y$ with
$\pi = \softmax(z)$, $\pi_k = e^{z_k}/\sum_je^{z_j}$. The loss is the categorical cross-entropy
\[
  \ell(z, y) = -\log\softmax(z)_y = \log\sum_{j=1}^Ke^{z_j} - z_y,\qquad
  \frac{\partial\ell}{\partial z_k} = \softmax(z)_k - \bbone[k = y].
\]
Both output units, and the derivations of these formulas, are the subject of
\cref{ch:sigmoid-softmax}.

\paragraph{Poisson.} For counts $y\in\{0, 1, 2, \dots\}$ (visits, arrivals, mutations),
$p(y\mid\lambda) = \lambda^ye^{-\lambda}/y!$ with a positive rate $\lambda = e^z$:
\[
  \ell(z, y) = e^z - yz + \log y!,\qquad
  \frac{\partial\ell}{\partial z} = e^z - y = \lambda - y,\qquad
  \frac{\partial^2\ell}{\partial z^2} = e^z>0 .
\]
The term $\log y!$ does not depend on the parameters and is dropped in practice.
With a linear model $z = w^\top x$ this is Poisson regression, and the loss is convex
in $w$.

\begin{example}[A Poisson step]\label{output-distributions:ex:poisson}
A shop has $y = 7$ visitors in an hour, and the model predicts $z = \log5$, i.e.\
$\lambda = 5$. Then $\ell = 5 - 7\log5 + \log7!\approx5 - 11.266 + 8.525 = 2.259$, and
$\partial\ell/\partial z = 5 - 7 = -2$: gradient descent increases $z$, hence the predicted rate. At
$\lambda = 7$ the gradient vanishes; the minimizer of the loss over $\lambda$ is the observed
count.
\end{example}

\section{Why the gradients look alike: the exponential family}\label{output-distributions:sec:expfam}

The gradients above are $(z - y)/\sigma^2$, $\sigma(z) - y$, $\softmax(z) - e_y$ and $\lambda - y$: in
each case, \emph{predicted mean minus observation}. This is not a coincidence.

\begin{definition}[Exponential family]\label{output-distributions:def:expfam}
A family of distributions is an \emph{exponential family} with \emph{natural
parameter} $\eta\in\R^m$ and \emph{sufficient statistic} $T(y)\in\R^m$ if
\[
  p(y\mid\eta) = h(y)\exp\bigl(\eta^\top T(y) - A(\eta)\bigr),\qquad
  A(\eta) = \log\int h(y)e^{\eta^\top T(y)}\,dy
\]
(a sum for discrete $y$). The \emph{log-partition function} $A$ makes $p$ integrate
to $1$.
\end{definition}

\begin{proposition}[Gradient of the negative log-likelihood]\label{output-distributions:prop:expfam}
For an exponential family, $\ell(\eta, y) = -\log p(y\mid\eta) = A(\eta) - \eta^\top T(y) - \log h(y)$, and
\[
  \nabla_\eta\ell = \nabla A(\eta) - T(y) = \E_\eta[T(Y)] - T(y),\qquad
  \nabla^2_\eta\ell = \nabla^2A(\eta) = \Cov_\eta\bigl(T(Y)\bigr)\succeq0 .
\]
In particular, $\ell$ is convex in the natural parameter.
\end{proposition}

\begin{proof}
Differentiate $A(\eta) = \log Z(\eta)$, $Z(\eta) = \int h(y)e^{\eta^\top T(y)}\,dy$, under the integral sign:
\[
  \nabla A = \frac{\nabla Z}{Z} = \int T(y)\,\frac{h(y)e^{\eta^\top T(y)}}{Z(\eta)}\,dy = \E_\eta[T(Y)],
\]
\[
  \nabla^2A = \frac{\nabla^2Z}Z - \frac{\nabla Z\,\nabla Z^\top}{Z^2} = \E_\eta[T(Y)T(Y)^\top] - \E_\eta[T(Y)]\,\E_\eta[T(Y)]^\top = \Cov_\eta(T(Y)).
\]
A covariance matrix is positive semidefinite.
\end{proof}

When the output unit is chosen so that the network output \emph{is} the natural
parameter, $\eta = z$, the gradient with respect to $z$ is $\E[T(Y)] - T(y)$ and the loss
is convex in $z$. Such a choice is called the \emph{canonical link} (in the language
of generalized linear models). The four examples are of this kind:
\begin{center}
\begin{tabular}{@{}lllll@{}}
\toprule
distribution & $T(y)$ & natural parameter $\eta$ & $A(\eta)$ & output unit (mean) \\
\midrule
Gaussian, $\sigma = 1$ & $y$ & $\mu$ & $\eta^2/2$ & identity \\
Bernoulli & $y$ & $\log\frac\pi{1-\pi}$ (logit) & $\log(1 + e^\eta)$ & sigmoid \\
Poisson & $y$ & $\log\lambda$ & $e^\eta$ & exponential \\
categorical & $e_y$ & logits $z$ & $\log\sum_ke^{z_k}$ & softmax \\
\bottomrule
\end{tabular}
\end{center}
(For the categorical family the logits are defined only up to an additive
constant, a point taken up in \cref{sigmoid-softmax:sec:softmax}.) The Gaussian with
learned variance is also an exponential family, with $T(y) = (y, y^2)$, but the
parametrization $(\mu,\log\sigma^2)$ used above is not the natural one, and its loss is not
convex in $(z_1, z_2)$ jointly.

\begin{keyidea}[Canonical links]
Pairing an exponential family with its canonical output unit gives a loss that is
convex in the network's last pre-activation, with gradient
``predicted mean minus observed statistic''. The gradient vanishes exactly when
the model's prediction matches the data on average --- the same moment-matching
condition as in \cref{cross-entropy:exr:moments}.
\end{keyidea}

\section{Multimodal targets: mixture density networks}\label{output-distributions:sec:mdn}

A single Gaussian per input cannot represent a target with several plausible
values for the same $x$, as in inverse problems (the joint angles of a robot arm
that reach a given position), or in the trajectory of a car that may turn left
or right. Fitted by maximum likelihood, a single Gaussian then does what forward
KL does in \cref{cross-entropy:ex:bimodal}: it predicts the average of the
modes, which may be an impossible value. A \emph{mixture density network} (Bishop,
\emph{Mixture Density Networks}, Aston University technical report NCRG/94/004, 1994) outputs the parameters of a mixture
\[
  p(y\mid x) = \sum_{k=1}^K\pi_k(x)\,\mathcal N\bigl(y;\mu_k(x), \sigma_k^2(x)\bigr),
\]
with $\pi = \softmax(z^\pi)$, $\mu_k = z^\mu_k$ and $\sigma_k = e^{z^\sigma_k}$. The loss is
\[
  \ell = -\log\sum_{k=1}^K\exp\Bigl(\log\pi_k + \log\mathcal N(y;\mu_k,\sigma_k^2)\Bigr),
\]
computed with the log-sum-exp trick of \cref{sigmoid-softmax:sec:numerics}. Its
gradient involves the posterior probability that $y$ came from component $k$,
\[
  r_k = \frac{\pi_k\,\mathcal N(y;\mu_k,\sigma_k^2)}{\sum_j\pi_j\,\mathcal N(y;\mu_j,\sigma_j^2)},\qquad
  \frac{\partial\ell}{\partial\mu_k} = r_k\,\frac{\mu_k - y}{\sigma_k^2},\qquad
  \frac{\partial\ell}{\partial z^\pi_k} = \pi_k - r_k,
\]
the \emph{responsibilities} of the EM algorithm (\cref{ch:em-algorithm}): each component is
pulled towards $y$ in proportion to how much it claims the point
(\cref{output-distributions:exr:5}). Mixture density networks inherit the
collapse of the variance described above, now per component.

\paragraph{Other targets.} The same recipe applies to any family: a Beta
distribution for targets in $(0,1)$, with positive shape parameters $\alpha = e^{z_1}$,
$\beta = e^{z_2}$; a von Mises distribution for angles (\cref{ch:flat-circle}); a
negative binomial for overdispersed counts. For \emph{structured} outputs ---
a sentence, an image, a graph --- the joint distribution of the components is
factorized by the chain rule, $p(y_1,\dots,y_T\mid x) = \prod_tp(y_t\mid y_{<t}, x)$, and each factor
is one of the distributions above, typically a categorical over a vocabulary.
The loss is then a sum of cross-entropies, one per position; this is the loss of
the language models of \cref{part:language}.

\begin{table}[ht]
\centering
\caption{Output distributions, output units and losses. In each row the gradient
with respect to the (first) output is ``prediction minus target'', scaled.}
\label{output-distributions:tab:summary}
\small\setlength{\tabcolsep}{4pt}
\begin{tabular}{@{}lllll@{}}
\toprule
target & distribution & output unit & loss (up to constants) & $\partial\ell/\partial z$ \\
\midrule
real & Gaussian, $\sigma$ fixed & $\mu = z$ & $(y - z)^2/2\sigma^2$ & $(z - y)/\sigma^2$ \\
real & Laplace, $b$ fixed & $\mu = z$ & $\abs{y - z}/b$ & $\sign(z - y)/b$ \\
real & Gaussian, $\sigma$ learned & $\mu = z_1$, $\sigma^2 = e^{z_2}$ & $\frac{z_2}2 + \frac{(y - z_1)^2}{2e^{z_2}}$ & $(z_1 - y)/\sigma^2$ \\
binary & Bernoulli & $\pi = \sigma(z)$ & $\log(1 + e^z) - yz$ & $\sigma(z) - y$ \\
class & categorical & $\pi = \softmax(z)$ & $\log\sum_ke^{z_k} - z_y$ & $\softmax(z) - e_y$ \\
count & Poisson & $\lambda = e^z$ & $e^z - yz$ & $e^z - y$ \\
multimodal & mixture & softmax, id., exp & $-\log\sum_k\pi_k\mathcal N(y;\mu_k,\sigma_k^2)$ & $r_k(\mu_k - y)/\sigma_k^2$ \\
\bottomrule
\end{tabular}
\end{table}

\begin{chaptersummary}
\begin{itemize}[nosep]
  \item Designing a loss is choosing a conditional distribution for the target and
        an output unit that maps the network's outputs to valid parameters; the loss
        is the negative log-likelihood.
  \item Gaussian noise gives the squared error, Laplace noise the absolute error,
        Bernoulli and categorical models the binary and categorical cross-entropies,
        Poisson counts the Poisson loss.
  \item Learning the variance lets the network report its uncertainty, but the
        per-example loss is unbounded below: guard against collapse.
  \item For an exponential family, the negative log-likelihood is convex in the natural
        parameter with gradient $\E_\eta[T(Y)] - T(y)$. With the canonical link this is the
        familiar ``prediction minus target''.
  \item Multimodal targets need a mixture; its gradient weights each component by
        its responsibility.
\end{itemize}
\end{chaptersummary}

\section*{Exercises}
\addcontentsline{toc}{section}{Exercises}

\begin{exercise}\label{output-distributions:exr:1}
For the model $y_i\sim\mathcal N(\mu(x_i), \sigma^2)$ with a single $\sigma$, show that for fixed predictions
$\mu(x_i)$ the log-likelihood is maximized by $\hat\sigma^2 = \frac1n\sum_i(y_i - \mu(x_i))^2$, and compute the maximal
value of $\frac1n\log\mathcal L$ as a function of $\hat\sigma^2$.
\solutionlink{output-distributions:exr:1}
\end{exercise}

\begin{exercise}\label{output-distributions:exr:2}
For the Laplace model with a constant prediction $\mu$ for all points, show that
$\mu\mapsto\sum_i\abs{y_i - \mu}$ is minimized by any median of $y_1,\dots,y_n$. For the data
$1, 2, 3, 4, 100$, compare the maximum-likelihood constant under the Gaussian and the
Laplace models.
\solutionlink{output-distributions:exr:2}
\end{exercise}

\begin{exercise}\label{output-distributions:exr:3}
Write the Bernoulli and Poisson distributions in the form of
\cref{output-distributions:def:expfam}, identify $\eta$, $T$, $h$ and $A$, and check
\cref{output-distributions:prop:expfam} by computing $A'$ and $A''$ directly.
\solutionlink{output-distributions:exr:3}
\end{exercise}

\begin{exercise}\label{output-distributions:exr:4}
Show that for the Gaussian with learned variance, with $z_1$ fixed, the loss is convex
in $z_2 = \log\sigma^2$, but that the loss is not jointly convex in $(z_1, z_2)$.
\hint{Compute the Hessian at a point where $y\neq z_1$.}
\solutionlink{output-distributions:exr:4}
\end{exercise}

\begin{exercise}\label{output-distributions:exr:5}
Derive the gradients of the mixture density loss with respect to $\mu_k$ and to the
mixture logits $z^\pi_k$ given in \cref{output-distributions:sec:mdn}.
\solutionlink{output-distributions:exr:5}
\end{exercise}
\end{notechapter}
